\chapter{近似と広義積分}\label{章:近似と広義積分}

これまでは，積分を計算できる形へ変えることを中心に考えてきた。ところが，原始関数を知っている関数だけが積分の対象になるわけではない。また，原始関数が分かっていても，その値を小数で求めるには別の計算が必要になる。この章では，積分を「正確な値を表す式」として使うだけでなく，「近似したときに残る量を表す式」として使う。

もう一つ，定積分の定義で有限にしていた区間を広げる。上端をどこまでも遠ざけたり，被積分関数が限りなく大きくなる端点へ近づけたりするとき，途中までの積分が一定の値に近づくとは限らない。そこで必要になるのは，先に式を代入することではなく，どの極限を調べているかを明確にすることである。近似の誤差を評価することと，広義積分が収束するかを判定することは，どちらも積分の大小関係によって支えられている。

\section{有限の多項式で関数を表す}\label{節:有限テイラー}

\subsection{接線の次に何を加えるか}

微分係数による近似は
\[
 f(x)\simeq f(a)+f'(a)(x-a)
\]
という形であった。右辺は，点$a$で関数の値と傾きを合わせた一次式である。さらに二階微分の値も合わせたいなら，$(x-a)^2$の項を加える。$c(x-a)^2$を二回微分すると$2c$になるので，その係数は$f''(a)$そのものではなく，$f''(a)$の半分になる。同じ考えを続けると，$n$回までの微分の値を合わせた多項式は
\[
 P_n(x)=\sum_{k=0}^{n}\frac{f^{(k)}(a)}{k!}(x-a)^k
\]
となる。ここで$f^{(0)}=f$，$0!=1$とする。

この多項式が関数にどの程度近いかは，微分の値を合わせただけでは分からない。必要なのは，差$f(x)-P_n(x)$を調べる式である。まず微積分の基本定理から
\[
 f(x)=f(a)+\int_a^x f'(t)\dd t
\]
と書ける。積分の中で$f'(t)$を$f'(a)$に置き換えれば接線の式が現れるが，その置き換えで失われる分も残しておきたい。そのために，次の部分積分を行う。
\[
\begin{aligned}
 \int_a^x (x-t)f''(t)\dd t
 &=\bigl[(x-t)f'(t)\bigr]_a^x
   +\int_a^x f'(t)\dd t\\
 &=-(x-a)f'(a)+f(x)-f(a).
\end{aligned}
\]
したがって
\[
 f(x)=f(a)+f'(a)(x-a)
      +\int_a^x(x-t)f''(t)\dd t.
\]
最後の積分は，接線近似を正確な等式にするための項である。二階微分が小さい区間ではこの積分も小さくなる，という関係がここに現れている。

\begin{theorem}{有限テイラー公式と積分剰余}
$n$を$0$以上の整数とし，$f$は$a$と$x$を含む区間で$n+1$回連続微分可能であるとする。このとき
\[
 f(x)=\sum_{k=0}^{n}\frac{f^{(k)}(a)}{k!}(x-a)^k+R_n(x),
 \qquad
 R_n(x)=\frac{1}{n!}\int_a^x(x-t)^n f^{(n+1)}(t)\dd t
\]
が成り立つ。$R_n(x)$を剰余項という。
\end{theorem}

公式の仕組みを確かめる。$n=0$は微積分の基本定理そのものである。$n\geq1$のとき，剰余項を部分積分すると
\[
\begin{aligned}
 R_n(x)
 &=\left[\frac{(x-t)^n}{n!}f^{(n)}(t)\right]_a^x
   +\frac{1}{(n-1)!}\int_a^x(x-t)^{n-1}f^{(n)}(t)\dd t\\
 &=-\frac{f^{(n)}(a)}{n!}(x-a)^n+R_{n-1}(x).
\end{aligned}
\]
つまり，一つ前の剰余項から$n$次の項を取り出した残りが$R_n$である。この関係を順に使えば，公式が得られる。$x<a$でも，向きのある定積分の部分積分として同じ計算が成り立つ。

この導出で部分積分を選んだ理由も見ておこう。$(x-t)^n$を$t$で微分すれば次数が一つ下がり，$f^{(n+1)}(t)$を積分すれば微分の回数が一つ減る。その結果，一つ前の剰余項と結び付く。形式だけを覚えるよりも，この二つが同時に一段ずつ戻ることを覚えておくとよい。

%
\subsection{誤差の大きさと向き}

$a$と$x$の間で$|f^{(n+1)}(t)|\leq M$なら，積分の絶対値の評価から
\[
 |R_n(x)|\leq\frac{M}{n!}\int_{\min(a,x)}^{\max(a,x)}|x-t|^n\dd t
 =\frac{M|x-a|^{n+1}}{(n+1)!}
\]
を得る。多項式をどこまで計算すれば十分かは，この上限と，求められた精度を比較して決められる。$x$が$a$に近いほど高い次数の項が小さくなるだけでなく，最後に残る項も具体的な不等式で押さえられる。

なお，$x>a$では$(x-t)^n\geq0$である。したがって$f^{(n+1)}$の符号が分かれば，剰余項の符号も分かる。誤差の絶対値だけでなく，近似値が真の値より大きいか小さいかを判断できることがある。

\begin{example}{対数を三次式で近似する}
$0\leq x\leq\frac12$の範囲で$\log(1+x)$を三次式で近似し，誤差を評価せよ。
\end{example}

中心を$a=0$に選ぶ。$\log(1+x)$は$x=0$で値が$0$になり，各階の微分の値も簡単だからである。微分すると
\[
 f'(x)=\frac{1}{1+x},\quad
 f''(x)=-\frac{1}{(1+x)^2},\quad
 f'''(x)=\frac{2}{(1+x)^3},\quad
 f^{(4)}(x)=-\frac{6}{(1+x)^4}
\]
なので
\[
 \log(1+x)=x-\frac{x^2}{2}+\frac{x^3}{3}+R_3(x).
\]
$0\leq t\leq x$では$-6\leq f^{(4)}(t)<0$である。したがって
\[
 -\frac{x^4}{4}\leq R_3(x)\leq0,
\]
すなわち
\[
 x-\frac{x^2}{2}+\frac{x^3}{3}-\frac{x^4}{4}
 \leq\log(1+x)
 \leq x-\frac{x^2}{2}+\frac{x^3}{3}
\]
となる。例えば$x=\frac15$なら
\[
 \frac{137}{750}-\frac{1}{2500}
 \leq\log\frac65\leq\frac{137}{750}.
\]
三次式の値を計算してから，残りの最大値を差し引けば，真の値を含む区間が得られる。

\begin{example}{原始関数の公式を使わずに積分を近似する}
$A=\int_0^1\e^{-x^2}\dd x$を有理数で上下から評価せよ。
\end{example}

$\e^{-x^2}$を直接何回も微分すると，多項式と指数関数の積が増えていく。ここでは，まず一変数の関数$\e^{-u}$を$u=0$で展開し，最後に$u=x^2$を代入する。$0\leq u\leq1$に対して
\[
 \e^{-u}=1-u+\frac{u^2}{2}-\frac{u^3}{6}+r_3(u).
\]
四階微分は$\e^{-u}$であり，この区間で$0<\e^{-u}\leq1$だから
\[
 0\leq r_3(u)\leq\frac{u^4}{24}.
\]
$u=x^2$を代入して積分すると
\[
\begin{aligned}
 A&=\int_0^1\left(1-x^2+\frac{x^4}{2}-\frac{x^6}{6}\right)\dd x
       +\int_0^1r_3(x^2)\dd x\\
  &=1-\frac13+\frac1{10}-\frac1{42}
       +\int_0^1r_3(x^2)\dd x.
\end{aligned}
\]
ここで
\[
 0\leq\int_0^1r_3(x^2)\dd x
 \leq\frac1{24}\int_0^1x^8\dd x=\frac1{216}
\]
なので
\[
 \boxed{\frac{26}{35}\leq A\leq\frac{26}{35}+\frac1{216}}.
\]
積分区間全体で成立する誤差評価を用意したので，その不等式を積分できた。何点かで近似がよいことを確かめるだけでは，この結論にはならない。

%
\subsection{有限の公式と無限級数を分ける}

有限テイラー公式は，剰余項を含めた等式である。これに対して
\[
 f(x)=\sum_{k=0}^{\infty}\frac{f^{(k)}(a)}{k!}(x-a)^k
\]
と書くには，$n\to\infty$のとき$R_n(x)\to0$になることを確かめる必要がある。有限の段階で正しい公式を何度でも書けることと，その極限で剰余項が消えることは別の主張である。

この違いは，等比数列の和にすでに現れている。$x\ne1$に対して
\[
 \frac1{1-x}=1+x+\cdots+x^n+\frac{x^{n+1}}{1-x}
\]
は有限の恒等式である。$|x|<1$なら最後の項は$0$に近づくので，無限級数の和も$\frac1{1-x}$になる。しかし$x=2$なら，最後の項は$-2^{n+1}$であり，小さくならない。有限の恒等式は正しいままであっても，剰余項を捨てることはできない。

テイラー公式で$|f^{(n+1)}(t)|\leq M_n$と評価したときも，同じことを調べる。$|x-a|^{n+1}$だけを見るのではなく，
\[
 \frac{M_n|x-a|^{n+1}}{(n+1)!}
\]
全体が$0$へ近づくかを調べなければならない。微分を繰り返すにつれて$M_n$が大きくなる場合もあるからである。

例えば$0\leq x\leq1$での$\log(1+x)$については，$(n+1)$階微分の絶対値が$n!$以下なので
\[
 |R_n(x)|\leq\frac{x^{n+1}}{n+1}\leq\frac1{n+1}.
\]
よって，この区間では剰余項が$0$に近づき，$x=1$も含めてテイラー級数が元の関数を表すことを証明できる。

\begin{note}
何回でも微分できる関数が，必ずテイラー級数と一致するわけでもない。例えば
\[
 h(x)=
 \begin{cases}
  \e^{-\frac1{x^2}}&(x\ne0),\\
  0&(x=0)
 \end{cases}
\]
は，$x=0$でのすべての階の微分係数が$0$になる。それでも$x\ne0$では$h(x)>0$である。そのため，$0$を中心に作ったテイラー多項式はすべて$0$であり，元の関数と一致しない。

この性質の理由は，指数関数の減少が任意のべきより速いことにある。$z>0$では，有限テイラー公式の剰余項の正値性から$\e^z\geq\frac{z^N}{N!}$である。$N$を十分大きく選べば，任意の固定した$m$について
\[
 \lim_{x\to0}|x|^{-m}\e^{-\frac1{x^2}}=0
\]
が従う。$x\ne0$で$h$を繰り返し微分すると，$x$の負の整数べきの有限和に$\e^{-\frac1{x^2}}$を掛けた形になる。この極限を，各階の微分係数の定義に順に使うと，$x=0$でどの階の微分係数も$0$になることが分かる。無限に滑らかであることと，テイラー級数で表せることは区別されるのである。
\end{note}

%
%
\section{台形公式とシンプソン公式}\label{節:数値積分}

\subsection{曲線を直線に置き換える}

定積分を近似する最初の方法は，区間を細かくして長方形の面積を足すことであった。両端の関数値を使えば，長方形の代わりに台形を使える。区間$[a,b]$で，両端の点を通る直線は
\[
 L(x)=\frac{b-x}{b-a}f(a)+\frac{x-a}{b-a}f(b)
\]
である。この直線の下の符号付き面積は
\[
 T=\int_a^bL(x)\dd x=\frac{b-a}{2}\{f(a)+f(b)\}
\]
となる。これが一つの区間に対する台形公式である。

図\ref{図:台形近似}では，曲線が下に凸なので弦が曲線より上にある。この場合，台形公式の値は積分の値より大きくなる。ただし，関数によって凸凹が変わる場合もある。そこで，図だけでなく二階微分による誤差の式を求めよう。以下の台形公式の誤差評価では，$f$ は $[a,b]$ で二回連続微分可能であるとする。

\begin{figure}[htbp]
\centering
\begin{tikzpicture}[x=1.7cm,y=1.4cm]
 \fill[blue!10] (0.5,0)--(2.5,0)--(2.5,1.9625)--(0.5,0.4625)--cycle;
 \draw[-{Latex}] (-0.1,0)--(2.9,0) node[right] {$x$};
 \draw[-{Latex}] (0,-0.1)--(0,2.35) node[above] {$y$};
 \draw[dashed] (0.5,0)--(0.5,0.4625);
 \draw[dashed] (2.5,0)--(2.5,1.9625);
 \draw[thick,orange!80!black] (0.5,0.4625)--(2.5,1.9625);
 \draw[thick,blue!70!black,domain=0.15:2.65,samples=80] plot (\x,{0.4+0.25*\x*\x});
 \fill (0.5,0.4625) circle (1.5pt);
 \fill (2.5,1.9625) circle (1.5pt);
 \node[below] at (0.5,0) {$a$};
 \node[below] at (2.5,0) {$b$};
 \node[above left] at (1.55,1.25) {$L(x)$};
 \node[right] at (2.05,1.18) {$f(x)$};
\end{tikzpicture}
\caption{端点を結ぶ直線の積分が台形公式になる。下に凸の曲線では，弦と曲線の間の面積だけ過大になる。}
\label{図:台形近似}
\end{figure}

$w(x)=(x-a)(b-x)$と置く。端点で$w$が$0$となり，$w''(x)=-2$となることを利用して，二回部分積分する。
\[
\begin{aligned}
 \int_a^b w(x)f''(x)\dd x
 &=-\int_a^bw'(x)f'(x)\dd x\\
 &=-[w'(x)f(x)]_a^b+\int_a^bw''(x)f(x)\dd x\\
 &=(b-a)\{f(a)+f(b)\}-2\int_a^bf(x)\dd x.
\end{aligned}
\]
したがって，誤差を「積分の値から近似値を引いたもの」と決めれば
\[
 \boxed{\int_a^bf(x)\dd x-T
 =-\frac12\int_a^b(x-a)(b-x)f''(x)\dd x}.
\]
$f''\geq0$なら右辺は$0$以下になるので，図で見た大小関係が式からも分かる。また$|f''(x)|\leq M_2$なら
\[
 \left|\int_a^bf(x)\dd x-T\right|
 \leq\frac{M_2}{2}\int_a^b(x-a)(b-x)\dd x
 =\frac{M_2(b-a)^3}{12}.
\]

区間全体を$N$等分し，$h=\frac{b-a}{N}$，$x_j=a+jh$とすると，各小区間の台形を足した値は
\[
 T_N=h\left\{\frac{f(a)+f(b)}2+\sum_{j=1}^{N-1}f(x_j)\right\}.
\]
内部の点は隣り合う二つの台形で半分ずつ使われるため，係数が$1$になる。各小区間の誤差は$\frac{M_2h^3}{12}$以下なので，合計では
\[
 \left|\int_a^bf(x)\dd x-T_N\right|
 \leq\frac{M_2(b-a)h^2}{12}
 =\frac{M_2(b-a)^3}{12N^2}.
\]
分割数を二倍にすれば，この誤差上限は四分の一になる。

%
\subsection{三点を放物線で結ぶ}

次に，区間の両端に加えて中点の値も使う。中点を$m$，半分の幅を$h$とし，区間を$[m-h,m+h]$と書く。$y=x-m$によって中心を$0$へ移して，三点を通る二次式を
\[
 q(y)=A+By+Cy^2
\]
とする。$q(0)=f(m)$より$A=f(m)$である。また，$q(h)$と$q(-h)$を足せば$B$の項が消えるので
\[
 C=\frac{f(m-h)-2f(m)+f(m+h)}{2h^2}.
\]
積分でも$By$は奇関数として消える。したがって
\[
\begin{aligned}
 \int_{-h}^hq(y)\dd y
 &=2hA+\frac{2h^3}{3}C\\
 &=\frac h3\{f(m-h)+4f(m)+f(m+h)\}.
\end{aligned}
\]
これがシンプソン公式である。中点の係数$4$は，二次式の積分を計算した結果として決まる。

この公式は三次式についても正確である。中心を移した三次式のうち，三次と一次の項は奇関数なので，積分も両端と中点の重み付き和も$0$になる。残った二次と定数の項については，すでに正確である。この「三次式まで誤差がない」という性質により，誤差は四階微分で評価できる。

%
\subsection{シンプソン公式の誤差を導く}

単に公式の係数を覚えるだけでは，どれだけ細かく分ければよいかは決められない。ここでは，有限テイラー公式から誤差の定数も導く。座標を移して区間を$[-h,h]$とし，$f$は四回連続微分可能とする。$-h$での三次テイラー多項式を$P$とすると
\[
 f(x)=P(x)+\frac16\int_{-h}^x(x-t)^3f^{(4)}(t)\dd t.
\]
また，$f$の原始関数を$F$，$-h$での$F$の四次テイラー多項式を$Q$とする。$F'=f$だから$Q'=P$であり，$F^{(5)}=f^{(4)}$である。$F$の四次の有限テイラー公式を$x=h$で使い，
\[
 F(h)-F(-h)=\int_{-h}^h f(x)\dd x,\qquad
 Q(h)-Q(-h)=\int_{-h}^h P(x)\dd x
\]
を合わせると
\[
 \int_{-h}^hf(x)\dd x
 =\int_{-h}^hP(x)\dd x
  +\frac1{24}\int_{-h}^h(h-t)^4f^{(4)}(t)\dd t
\]
を得る。この方法なら，二重積分の積分順序を入れ替える定理を新しく使う必要はない。

一方，最初の式を$x=-h,0,h$で使い，シンプソン公式の重みを掛けて足す。$P$に対する誤差は$0$なので，残るのは
\[
 \int_{-h}^hf(x)\dd x
 -\frac h3\{f(-h)+4f(0)+f(h)\}
 =\int_{-h}^hK(t)f^{(4)}(t)\dd t
\]
である。ただし，係数をまとめた$K$は
\[
 K(t)=
 \begin{cases}
 -\frac{(h+t)^3(h-3t)}{72}&(-h\leq t\leq0),\\[4pt]
 -\frac{(h-t)^3(h+3t)}{72}&(0\leq t\leq h).
 \end{cases}
\]
例えば$t\geq0$では，中点$f(0)$の剰余項に$t$は現れないので，係数は
\[
 \frac{(h-t)^4}{24}-\frac h{18}(h-t)^3
 =-\frac{(h-t)^3(h+3t)}{72}
\]
となる。$t\leq0$では，これに中点からの項$-\frac{2h}{9}(-t)^3$を加え，整理すると上の式になる。

$K$は偶関数であり，全区間で$K\leq0$である。その積分の絶対値は
\[
\begin{aligned}
 \int_{-h}^h|K(t)|\dd t
 &=\frac1{36}\int_0^h(h-t)^3(h+3t)\dd t\\
 &=\frac{h^5}{36}\int_0^1(1-u)^3(1+3u)\dd u
 =\frac{h^5}{90}.
\end{aligned}
\]
したがって$|f^{(4)}|\leq M_4$なら，一組の区間について
\[
 \left|\int_{m-h}^{m+h}f(x)\dd x
 -\frac h3\{f(m-h)+4f(m)+f(m+h)\}\right|
 \leq\frac{M_4h^5}{90}
\]
となる。さらに$f^{(4)}\geq0$なら，積分の値から近似値を引いた誤差は$0$以下である。

\begin{theorem}{合成シンプソン公式と誤差}
$N$を正の偶数とし，$h=\frac{b-a}{N}$，$x_j=a+jh$とする。$f$が$[a,b]$で四回連続微分可能で，$|f^{(4)}(x)|\leq M_4$なら
\[
 S_N=\frac h3\left\{f(a)+f(b)
 +4\sum_{j=1}^{\frac N2}f(x_{2j-1})
 +2\sum_{j=1}^{\frac N2-1}f(x_{2j})\right\}
\]
に対して
\[
 \left|\int_a^bf(x)\dd x-S_N\right|
 \leq\frac{M_4(b-a)h^4}{180}
 =\frac{M_4(b-a)^5}{180N^4}
\]
が成り立つ。
\end{theorem}

小区間を二つずつ組にすると組数は$\frac N2$である。各組の誤差上限$\frac{M_4h^5}{90}$を足すことで，この評価が得られる。組どうしの境目は左右の組で一回ずつ使われるため係数が$2$，各組の中点は係数が$4$になる。分割数を二倍にすると誤差上限は十六分の一になるが，これは四階微分が区間全体で有界という条件の下での結論である。端点で微分が発散する関数に，そのままこの上限を使うことはできない。

\begin{example}{同じ点の値から二つの近似を作る}
$\int_0^1\frac1{1+x}\dd x$について，四等分した台形公式とシンプソン公式を計算し，誤差を評価せよ。
\end{example}

分点は$0,\frac14,\frac12,\frac34,1$で，関数値は順に$1,\frac45,\frac23,\frac47,\frac12$である。したがって
\[
 T_4=\frac14\left(\frac34+\frac45+\frac23+\frac47\right)
 =\frac{1171}{1680},
\]
\[
 S_4=\frac1{12}\left\{\frac32+4\left(\frac45+\frac47\right)+\frac43\right\}
 =\frac{1747}{2520}.
\]
二階微分は$\frac2{(1+x)^3}$，四階微分は$\frac{24}{(1+x)^5}$で，いずれも正である。そのため，どちらの近似値も真の値以上になる。上限$M_2=2$，$M_4=24$を使えば
\[
 T_4-\frac1{96}\leq\int_0^1\frac1{1+x}\dd x\leq T_4,
\]
\[
 S_4-\frac1{1920}\leq\int_0^1\frac1{1+x}\dd x\leq S_4.
\]
この積分は$\log2$に等しいが，その小数値を知らなくても近似の保証は作れる。同じ五つの関数値を使っていても，曲線をどの次数の式で置き換えるかによって誤差の評価が変わる。

%
%
\section{商の極限とロピタルの定理}\label{節:ロピタル}

\subsection{微分した商を調べてよい条件}

極限を計算していると，分子も分母も$0$へ近づく商が現れる。その場合，分子と分母がどれだけ速く$0$へ近づくかを比べる必要がある。微分係数がその速さを表すことから，分子と分母の微分を比べる方法が考えられる。ただし，元の商の形を確かめず，機械的に分子と分母を微分してはいけない。

\begin{theorem}{ロピタルの定理の一側極限版}
$x$が区間の一端$c$へ近づく一側極限を考える。$c$は有限でも$+\infty,-\infty$でもよい。端点を除くその区間で$f,g$は微分可能であり，$g'(x)\ne0$とする。さらに，次のいずれかを満たすとする。
\begin{enumerate}[label=(\roman*),leftmargin=*]
 \item $f(x)\to0$かつ$g(x)\to0$である。
 \item $|f(x)|\to\infty$かつ$|g(x)|\to\infty$である。
\end{enumerate}
このとき
\[
 \lim_{x\to c}\frac{f'(x)}{g'(x)}=L
\]
が存在すれば
\[
 \lim_{x\to c}\frac{f(x)}{g(x)}=L
\]
である。ここで$L$は実数のほか，$+\infty,-\infty$でもよい。元の商は端点の十分近くで定義されている。
\end{theorem}

両側極限を求めるときは，左右の一側極限にそれぞれ適用し，結果が一致することを確認する。また，$\frac{f'}{g'}$をさらに微分して計算したい場合は，その新しい商についても同じ条件を改めて確認する。微分後の商が$0$と$0$の商になるとは限らない。

%
\subsection{コーシーの平均値の定理による根拠}

通常の平均値の定理では，関数の増分を区間の幅で割った。ここでは，二つの関数の増分を比べる。$f,g$が$[u,v]$で連続，内部で微分可能で，$g'\ne0$なら，ある$\xi\in(u,v)$が存在して
\[
 \frac{f(v)-f(u)}{g(v)-g(u)}=\frac{f'(\xi)}{g'(\xi)}
\]
となる。これをコーシーの平均値の定理という。

この式もロルの定理から導ける。関数
\[
 H(t)=\{f(v)-f(u)\}\{g(t)-g(u)\}
      -\{g(v)-g(u)\}\{f(t)-f(u)\}
\]
を作ると，$H(u)=H(v)=0$である。よって$H'(\xi)=0$となる点があり，微分して整理すると求める等式になる。なお，$g(v)=g(u)$であればロルの定理からどこかで$g'=0$となってしまうので，分母$g(v)-g(u)$は$0$でない。

まず有限の端点$c$で$0$と$0$の形になる場合を考える。端点に$f(c)=g(c)=0$を補えば連続になるので，$c$と$x$の間でコーシーの平均値の定理を使える。すると
\[
 \frac{f(x)}{g(x)}=\frac{f'(\xi_x)}{g'(\xi_x)}
\]
となり，$\xi_x$は$c$と$x$の間にある。$x\to c$なら$\xi_x\to c$だから，右辺の極限が$L$であることから結論が従う。

無限大の形では，少し見方を変える。例えば$x\to+\infty$で$|g(x)|\to\infty$，$L$が実数の場合，$\frac{f'(t)}{g'(t)}$が$L$に十分近い範囲の始点$B$を固定する。$x>B$に対して
\[
 D_x=\frac{f(x)-f(B)}{g(x)-g(B)}
\]
は，その範囲のどこかでの$\frac{f'}{g'}$に等しいので，$L$に十分近い。元の商は
\[
 \frac{f(x)}{g(x)}
 =D_x\left(1-\frac{g(B)}{g(x)}\right)+\frac{f(B)}{g(x)}
\]
と書ける。$B$を固定した後で$x\to+\infty$とすれば，固定した値を含む二つの比はそれぞれ$0$へ近づく。そのため元の商も$L$へ近づく。有限の端点で$|g|\to\infty$となる場合も，$B$を端点の十分近くに固定して同じ議論を使える。

$0$と$0$の形で端点が$+\infty$の場合は，十分大きい$x<y$の間に平均値の定理を使う。増分の比が$L$に近いことを確かめた後，$x$を固定して$y\to+\infty$とすると，その比は$\frac{f(x)}{g(x)}$へ近づく。$L=+\infty$の場合は「$L$との差が小さい」の代わりに「任意の大きな定数より大きい」を用い，$L=-\infty$では不等号を逆にする。同じ議論で拡張された極限も得られる。

\begin{example}{消える次数を比較する}
$\lim_{x\to0}\cfrac{\log(1+x)-x+\cfrac{x^2}{2}}{x^3}$を求めよ。
\end{example}

分子も分母も$0$に近づき，$x\ne0$で分母の微分$3x^2$は$0$でない。分子を微分すると
\[
 \frac1{1+x}-1+x=\frac{x^2}{1+x}
\]
とまとめられるので
\[
 \cfrac{\left(\log(1+x)-x+\cfrac{x^2}{2}\right)'}{(x^3)'}
 =\frac1{3(1+x)}\longrightarrow\frac13.
\]
よって求める極限は$\frac13$である。微分後に式を整理すると$x^2$が約分できるため，三回の微分は必要ない。

同じ極限は有限テイラー公式でも説明できる。三次の項まで書けば，分子は$\frac{x^3}{3}+R_3(x)$となる。$0$の近くでは$|R_3(x)|\leq C|x|^4$と評価できるので，$x^3$で割った剰余項は$0$へ近づく。どちらの方法も，分子の最初に残る項が三次であることを調べている。

ロピタルの定理の条件は省略できない。例えば$x\to0$で$\frac{x^2+1}{x^2}$を考えると，微分した商は$1$だが，元の商は$+\infty$へ発散する。元の分子は$0$にも無限大にも近づかず，定理の仮定を満たしていない。また$x\to+\infty$で
\[
 \frac{x+\sin x}{x}\longrightarrow1
\]
であっても，微分した商$1+\cos x$には極限がない。微分した商の極限が存在しないことから，元の商の極限が存在しないとは結論できない。

%
%
\section{端点を極限にする積分}\label{節:広義積分の定義}

\subsection{有限の途中まで積分してから極限をとる}

$\int_0^1\frac1{\sqrt{x}}\dd x$では，$x=0$で被積分関数が定義されず，その近くで値が限りなく大きくなる。それでも，$\varepsilon>0$で切り取った区間では
\[
 \int_{\varepsilon}^1\frac1{\sqrt{x}}\dd x
 =2-2\sqrt{\varepsilon}
\]
と計算できる。この値が$\varepsilon\to0+$で$2$へ近づくことを利用して，積分の意味を広げる。

\begin{definition}{広義積分}
$f$が$[a,b)$の各閉部分区間で積分可能で，端点$b$に特異性をもつとき，
\[
 \int_a^bf(x)\dd x=\lim_{c\to b-}\int_a^cf(x)\dd x
\]
と定める。ただし，右辺が有限の実数として存在するとき，この広義積分は収束するという。左端に特異性がある場合も同様に，一側極限で定める。

上端が無限大の場合は
\[
 \int_a^{\infty}f(x)\dd x=\lim_{R\to\infty}\int_a^Rf(x)\dd x
\]
と定める。下端が$-\infty$の場合も同様である。有限の極限が存在しなければ発散するという。
\end{definition}

ここで，$+\infty$へ発散する場合も，振動して極限がない場合も，「収束しない」という点では同じである。一方，発散の仕方を区別したいときには「$+\infty$へ発散する」などと具体的に述べる。

%
\subsection{複数の特異点は独立に調べる}

両端に特異性がある場合は，内部の一点$c$を固定して
\[
 \int_a^bf(x)\dd x
 =\int_a^cf(x)\dd x+\int_c^bf(x)\dd x
\]
と定める。ただし，右辺の二つの広義積分がそれぞれ収束することが必要である。両端を同じ速さで動かして相殺することは，この定義には含まれない。分割点$c$を別の内部の点へ変えても，その間の通常の定積分が一方に加わり，他方から引かれるだけなので，値は変わらない。

区間内部に特異点$d$があれば，$d$の左右にも分ける。また，実数全体の積分は，例えば
\[
 \int_{-\infty}^{\infty}f(x)\dd x
 =\int_{-\infty}^{0}f(x)\dd x+\int_0^{\infty}f(x)\dd x
\]
として，左右を独立に調べる。複数の特異点があるときは，それぞれに別の切り取り幅や上限を用意して考えると，何を確認すべきかが明確になる。

この違いを$\frac1x$で確かめよう。$\varepsilon>0$に対して
\[
 \int_{-1}^{-\varepsilon}\frac1x\dd x=\log\varepsilon\longrightarrow-\infty,
 \qquad
 \int_{\varepsilon}^{1}\frac1x\dd x=-\log\varepsilon\longrightarrow+\infty.
\]
したがって$\int_{-1}^1\frac1x\dd x$は広義積分として収束しない。しかし，同じ$\varepsilon$を使って足せば常に$0$になる。このような対称的な切り取り方による極限を，コーシーの主値という。
\[
 \operatorname{PV}\int_{-1}^1\frac1x\dd x
 =\lim_{\varepsilon\to0+}
 \left(\int_{-1}^{-\varepsilon}\frac1x\dd x
       +\int_{\varepsilon}^{1}\frac1x\dd x\right)=0.
\]
主値が存在しても，広義積分が収束するとは限らない。実際，左を$\varepsilon$，右を$2\varepsilon$で切れば，二つの積分の和は$-\log2$になる。切り取り方を結び付けて得られた相殺だけでは，各側の収束を示したことにならない。

無限区間でも同じ注意が要る。例えば$\frac{x}{1+x^2}$は奇関数なので，任意の$R$に対して$\int_{-R}^R\frac{x}{1+x^2}\dd x=0$である。一方，右半分の積分は$\frac12\log(1+R^2)$となり発散する。奇関数であることを使う前に，広義積分の収束を確認する必要がある。

%
%
\section{べき関数との比較と条件収束}\label{節:広義積分の判定}

\subsection{基準になる二つの積分}

無限大での収束を調べる基準は
\[
 \int_1^{\infty}\frac1{x^p}\dd x
\]
である。$p\ne1$なら，まず有限の$R$まで積分して
\[
 \int_1^R x^{-p}\dd x=\frac{R^{1-p}-1}{1-p}.
\]
$p>1$なら$R^{1-p}\to0$なので収束し，値は$\frac1{p-1}$となる。$p<1$なら$R^{1-p}\to\infty$なので発散する。$p=1$では$\log R\to\infty$であり，やはり発散する。

これに対して，有限端点$0$の近くでは
\[
 \int_{\varepsilon}^{1}x^{-p}\dd x
 =\frac{1-\varepsilon^{1-p}}{1-p}\qquad(p\ne1)
\]
である。今度は$\varepsilon\to0+$なので，$p<1$のときに収束し，値は$\frac1{1-p}$になる。$p\geq1$では発散する。

\begin{center}
\begin{tabular}{lll}
\toprule
積分 & 収束する条件 & 収束時の値\\
\midrule
$\int_1^{\infty}x^{-p}\dd x$ & $p>1$ & $\frac1{p-1}$\\[4pt]
$\int_0^1x^{-p}\dd x$ & $p<1$ & $\frac1{1-p}$\\
\bottomrule
\end{tabular}
\end{center}

同じべき関数でも，不等号の向きが逆になる。遠くで減る速さを調べる場合と，端点で大きくなる速さを調べる場合では，問題が異なるからである。覚えるときは，$\frac1x$がどちらでも境目になることと，実際の原始関数の極限を対応させたい。

%
\subsection{比較判定法}

$0\leq f(x)\leq g(x)$が端点の十分近くで成り立つとする。このとき，$g$の広義積分が収束すれば$f$の広義積分も収束する。逆向きには，$f$の広義積分が$+\infty$へ発散すれば，$g$の広義積分も$+\infty$へ発散する。

例えば無限区間では，$\int_A^R f(x)\dd x$は$R$とともに増加する。しかも
\[
 0\leq\int_A^Rf(x)\dd x\leq\int_A^Rg(x)\dd x
 \leq\int_A^{\infty}g(x)\dd x
\]
と一定の上限をもつ。上に有界な単調増加の量は有限の極限をもつので，$f$の積分も収束する。端点付近でだけ比較できればよいのは，それ以外の閉区間の積分は有限の一定値を加えるだけだからである。

特に$f,g>0$で
\[
 \lim\frac{f(x)}{g(x)}=c,\qquad0<c<\infty
\]
となるなら，端点の十分近くで$\frac c2g(x)\leq f(x)\leq\frac{3c}{2}g(x)$と挟める。そのため二つの積分は同時に収束するか，同時に発散する。この極限比較では，$c$が正の有限値であるという条件が大切である。比の極限が$0$の場合は，$f$が$g$より小さいことしか分からず，両者の収束が一致するとは限らない。

\begin{example}{収束を先に確かめてから置換する}
$\int_0^1\frac1{\sqrt{x(1+x)}}\dd x$の収束を確かめ，値を求めよ。
\end{example}

特異性があるのは$x=0$だけである。$0<x\leq1$で
\[
 0<\frac1{\sqrt{x(1+x)}}\leq\frac1{\sqrt{x}}
\]
となり，右辺は$0$の近くで積分可能だから，求める広義積分は収束する。

計算では，まず$\varepsilon\leq x\leq1$に限って$x=t^2$と置く。平方根の$x$と$\dd x=2t\dd t$が約分されるので
\[
 \int_{\varepsilon}^1\frac1{\sqrt{x(1+x)}}\dd x
 =2\int_{\sqrt{\varepsilon}}^1\frac1{\sqrt{1+t^2}}\dd t.
\]
右辺の被積分関数は$t=0$でも連続である。$\varepsilon\to0+$として
\[
 2\int_0^1\frac1{\sqrt{1+t^2}}\dd t
 =2\left[\log\left(t+\sqrt{1+t^2}\right)\right]_0^1
 =2\log(1+\sqrt2)
\]
を得る。置換はまず通常の定積分に行い，その後で端点の極限をとった。

%
\subsection{符号の変化が収束を生む場合}

関数が正負に変わると，絶対値を積分した場合と，元の関数を積分した場合とで，収束の結果が異なることがある。

\begin{definition}{絶対収束と条件収束}
広義積分$\int f(x)\dd x$について，$\int|f(x)|\dd x$が収束するとき絶対収束するという。$\int f(x)\dd x$は収束するが，$\int|f(x)|\dd x$は発散するとき，条件収束するという。
\end{definition}

絶対収束なら元の積分も収束する。例えば無限区間では，遠い二点$S<R$の間の積分が
\[
 \left|\int_S^Rf(x)\dd x\right|\leq\int_S^R|f(x)|\dd x
\]
で押さえられる。右辺は$S,R$が十分大きければいくらでも小さくなるので，途中までの積分はコーシーの収束条件を満たす。これは実数の完備性による収束判定であり，「途中までの値が互いにいくらでも近くなるなら，有限の極限をもつ」という事実を使っている。

\begin{example}{振動する積分の条件収束}
$\int_1^{\infty}\frac{\sin x}{x}\dd x$は条件収束することを示せ。
\end{example}

$\frac1x$だけでは無限大で積分可能でない。しかし，$\sin x$を積分すると有界な$-\cos x$になり，$\frac1x$を微分すると$-\frac1{x^2}$になる。この二つを利用するため，部分積分を選ぶ。
\[
\begin{aligned}
 \int_1^R\frac{\sin x}{x}\dd x
 &=\left[-\frac{\cos x}{x}\right]_1^R
   -\int_1^R\frac{\cos x}{x^2}\dd x\\
 &=\cos1-\frac{\cos R}{R}-\int_1^R\frac{\cos x}{x^2}\dd x.
\end{aligned}
\]
境界項$\frac{\cos R}{R}$は$0$へ近づく。また$\frac{|\cos x|}{x^2}\leq\frac1{x^2}$なので，最後の積分は絶対収束する。よって元の積分は収束する。

絶対値を付けた積分は別に調べる。整数$k\geq1$について，$[k\pi,(k+1)\pi]$では$x\leq(k+1)\pi$だから
\[
 \int_{k\pi}^{(k+1)\pi}\frac{|\sin x|}{x}\dd x
 \geq\frac1{(k+1)\pi}\int_{k\pi}^{(k+1)\pi}|\sin x|\dd x
 =\frac2{(k+1)\pi}.
\]
この下限を足すと調和級数の定数倍になり，発散する。したがって絶対値の積分は発散し，元の積分は条件収束である。ここでは積分の値そのものを求めていない。収束を証明することと，値を計算することは別の課題である。

%
%
\section{原始関数が存在することと式で書けること}\label{節:非初等原始関数}

$\e^{-x^2}$の原始関数を，これまでの置換や部分積分で求めようとしても，なじみのある有限個の式にはたどり着かない。だからといって，原始関数が存在しないわけではない。微積分の基本定理を使えば
\[
 F(x)=\int_0^x\e^{-t^2}\dd t
\]
はすべての実数で定義され，$F'(x)=\e^{-x^2}$となる。これ自体が一つの原始関数である。

「初等関数で表せない」というときは，数，変数，代数的演算，指数関数，対数関数，三角関数などを有限回組み合わせた通常の初等関数の範囲では表せない，という意味である。$\e^{-x^2}$には初等関数で表された原始関数がないことが知られている。ただし，その不可能性の証明には，本書の範囲を超える理論が必要である。ここではこの事実を採用し，置換をいくつか試して失敗したことを証明の代わりにはしない。

積分で定義した$F$についても，多くのことが分かる。被積分関数が正なので$F$は単調増加である。被積分関数が偶関数なので$F$は奇関数である。さらに
\[
 F''(x)=-2x\e^{-x^2}
\]
から，$x>0$では上に凸，$x<0$では下に凸である。有限区間での値はテイラー公式や数値積分で評価できるし，$x\geq1$では$\e^{-x^2}\leq\e^{-x}$なので，$\int_0^{\infty}\e^{-x^2}\dd x$の収束も比較によって分かる。

積分を含む定義は，計算を途中で諦めた記号ではない。微分，対称性，大小関係，近似，極限を調べるための出発点になる。原始関数の表示が難しいときには，求めたいものが「不定積分の式」なのか，「特定の区間での値」なのか，「誤差付きの近似」なのか，「収束するという事実」なのかを分けて考える。その区別により，使う道具も変わってくる。

%
\section*{章末問題}
\begin{enumerate}[label=\textbf{\thechapter.\arabic*},leftmargin=*]
 \item\label{問:十一対数近似}
 $P(x)=x-\frac{x^2}{2}+\frac{x^3}{3}$とする。$0\leq x\leq\frac12$で
 $0\leq P(x)-\log(1+x)\leq\frac{x^4}{4}$を示し，$\log\frac{11}{10}$を幅$\frac1{40000}$以下の有理数の区間に挟め。
 \item\label{問:十一指数積分近似}
 $\e^{-u}$の二次テイラー多項式を使い，$\int_0^{\frac12}\e^{-x^2}\dd x$の近似値と誤差の上限を求めよ。近似値が真の値の上下どちらにあるかも述べよ。
 \item\label{問:十一数値積分}
 $\int_0^2x^4\dd x$を四等分した台形公式とシンプソン公式で近似せよ。真の値との差を計算し，二階微分，四階微分による誤差上限と比較せよ。
 \item\label{問:十一ロピタル}
 次の極限をロピタルの定理で求め，適用条件を述べよ。
 \[
 \text{(a)}\ \lim_{x\to0}\frac{\e^x-1-x}{x^2},
 \qquad
 \text{(b)}\ \lim_{x\to\infty}\frac{\log x}{x^{\alpha}}\quad(\alpha>0).
 \]
 \item\label{問:十一両端比較}
 実数$p,q$に対し，$\int_0^1\frac1{x^p(1-x)^q}\dd x$が収束するための必要十分条件を求めよ。両端を独立に調べること。
 \item\label{問:十一対数尾部}
 実数$q$に対し，$\int_{\e}^{\infty}\frac1{x(\log x)^q}\dd x$の収束条件を求め，収束する場合の値を求めよ。
 \item\label{問:十一主値区別}
 $\int_{-1}^1\frac1x\dd x$と$\int_{-1}^1\frac{1-\cos x}{x}\dd x$について，広義積分としての収束と値を調べよ。前者の対称的な主値も求め，後者との違いを説明せよ。
 \item\label{問:十一積分定義関数}
 $F(x)=\int_0^x\e^{-t^2}\dd t$とする。$F$の奇偶，増減，凹凸を調べよ。また$x\geq0$で$x-\frac{x^3}{3}\leq F(x)\leq x$を示し，$\lim_{x\to\infty}F(x)$が有限であることを証明せよ。
\end{enumerate}

\smallskip\noindent 解答は\pageref{解答:第11章}ページ。
%
